% jincai/jincai.tex —— 进才中学高一上周末作业 1（试卷版/答案版）
% 试卷版：xelatex jincai.tex
% 答案版：xelatex "\def\isanswer{1}\input{jincai.tex}"

\documentclass[a4paper,11pt]{article}
\usepackage{common}

\begin{document}

\examtitlest{2026--2027 学年度进才中学高一上周末作业 1}{集合与常用逻辑用语}

\bigsec{一、填空题}

\begin{enumerate}
\setcounter{enumi}{2}

\item 方程组 $\begin{cases}2x-y=3 \\ x+2y=4\end{cases}$ 的解集是 \blankline[2em].

\ans{$\{(2,1)\}$}

\item 已知集合 $A=\left\{x\,\middle|\,\dfrac{6}{5-x}\in\Z,\,x\in\mathbb{N}\right\}$，试用列举法表示集合 $A$，则 $A=$ \blankline[3em].

\ans{$\{2,3,4,6,7,8,11\}$}

\item 已知 $U=\R$，$A=\{x\mid x>0\}$，$B=\{x\mid x\leqslant -1\}$，则：$\left(A\cap\overline{B}\right)\cup\left(B\cap\overline{A}\right)=$ \blankline[3em].

\ans{$\{x\mid x\leqslant -1\text{ 或 }x>0\}$}

\item 全集 $U=\{x\mid x\text{ 是不大于 }20\text{ 的素数}\}$，若 $A\cap\overline{B}=\{3,5\}$，$\overline{A}\cap B=\{7,19\}$，$\overline{A\cup B}=\{2,17\}$，则集合 $A=$ \blankline[3em].

\ans{$\{3,5,11,13\}$}

\ifanswer
\solbegin
\step{全集 $U=\{x\mid x\text{ 取不大于 }20\text{ 的质数}\}=\{2,3,5,7,11,13,17,19\}$，$A\cap\overline{B}=\{3,5\}$，$\overline{A}\cap B=\{7,19\}$，$\overline{A\cup B}=\{2,17\}$，由韦恩图得 $A=\{3,5,11,13\}$，$B=\{7,11,13,19\}$.}
\solend

\fi
\item 已知集合 $A=\{1,2,3,4,5,6\}$，$B=\{(x,y)\mid x\in A,\,y\in A,\,x-y\in A\}$，则 $B$ 中所含元素的个数为 \blankline[2em].

\ans{$15$}

\ifanswer
\solbegin
\step{$B=\{(2,1),\allowbreak(3,1),\allowbreak(3,2),\allowbreak(4,1),\allowbreak(4,2),\allowbreak(4,3),\allowbreak(5,1),\allowbreak(5,2),\allowbreak(5,3),\allowbreak(5,4),\allowbreak(6,1),\allowbreak(6,2),\allowbreak(6,3),\allowbreak(6,4),\allowbreak(6,5)\}$，共 15 个元素.}
\solend

\fi
\item 已知集合 $A=\left\{(x,y)\,\middle|\,\dfrac{y-3}{x-2}=1\right\}$，$B=\{(x,y)\mid y=ax+2\}$，且满足 $A\cap B=\varnothing$，则实数 $a$ 的值是 \blankline[2em].

\ans{$a=1\text{ 或 }a=\dfrac{1}{2}$}

\ifanswer
\solbegin
\step{当 $A=\varnothing$ 时，$y=ax+2$ 过 $(2,3)$，解得 $a=\dfrac{1}{2}$；}
\step{当 $A\neq\varnothing$ 时，$y=x+1$ 与 $y=ax+2$ 平行，解得 $a=1$.}
\step{综上，$a=1$ 或 $a=\dfrac{1}{2}$.}
\solend

\fi
\item $Q$ 是有理数集，集合 $M=\{x\mid x=a+\sqrt{3}b,\,a,b\in Q,\,x\neq 0\}$，在下列集合中：

① $\{x\mid x=\sqrt{3}t,\,t\in M\}$；\quad ② $\{x\mid x=t^2,\,t\in M\}$；

③ $\{x\mid x=x_1+x_2,\,x_1\in M,\,x_2\in M\}$；\quad ④ $\{x\mid x=x_1x_2,\,x_1\in M,\,x_2\in M\}$.

与集合 $M$ 相等的集合序号是 \blankline[2em].

\ans{④}

\ifanswer
\solbegin
\step{对于①，因为 $t\in M$，设 $t=a+\sqrt{3}b,\,a,b\in Q$，所以 $\sqrt{3}t=\sqrt{3}(a+\sqrt{3}b)=\sqrt{3}a+3\sqrt{3}b$，不妨取 $a=2,\,b=0$，则 $2\in M$，而 $\sqrt{6}\in\{x\mid x=\sqrt{3}t,\,t\in Q\}$，显然 $\sqrt{6}\notin M$，所以①与集合 $M$ 不相等；}
\step{对于②，令 $t=a+\sqrt{3}b,\,a,b\in Q$，则 $t^2=(a+\sqrt{3}b)^2=a^2+3b^2+2\sqrt{3}ab$，显然 $-1+\sqrt{3}\in M$，但 $-1+\sqrt{3}\in\{x\mid x=t^2,\,t\in M\}$，即②与集合 $M$ 不相等；}
\step{对于③，当 $x_1=a+\sqrt{3}b,\,x_2=-a-\sqrt{3}b,\,a,b\in Q$ 时，此时 $x=x_1+x_2=0$，即 $0\in\{x\mid x=x_1+x_2,\,x_1\in M,\,x_2\in M\}$，而集合 $M$ 中不包含元素 $0$，所以③与集合 $M$ 不相等；}
\step{对于④，令 $x_1=a_1+\sqrt{3}b_1\,(a_1,b_1\in Q,x_1\neq 0),\,x_2=a_2+\sqrt{3}b_2\,(a_2,b_2\in Q,x_2\neq 0)$，则 $x=x_1x_2=(a_1a_2+3b_1b_2)+\sqrt{3}(a_1b_2+a_2b_1)$，其中 $a_1a_2+3b_1b_2,\,a_1b_2+a_2b_1\in Q,\,x\neq 0$，所以④与集合 $M$ 相等；}
\step{故与集合 $M$ 相等的集合序号是④.}
\solend

\fi
\item 对于区间 $[a,b]$ 我们规定 $b-a$ 是这个区间的“长度”. 已知 $A$、$B$ 都是集合 $\{x\mid -1\leqslant x\leqslant 1\}$ 的子集，$A=\{x\mid a\leqslant x\leqslant a+1\}$，$B=\left\{x\,\middle|\,b-\dfrac{3}{2}\leqslant x\leqslant b\right\}$，则集合 $A\cap B$“长度”的取值范围是 \blankline[3em].

\ans{$\left[\dfrac{1}{2},1\right]$}

\ifanswer
\solbegin
\step{区间 $A$ 长度为 1，区间 $B$ 长度为 $\dfrac{3}{2}$.}
\step{当 $a=-1,\,b=\dfrac{1}{2}$ 时，$A\cap B=A=[-1,0]$，区间长度为 1，取得最大值；}
\step{当 $a=-1,\,b=1$ 时，$A\cap B=\left[-\dfrac{1}{2},0\right]$ 区间长度为 $\dfrac{1}{2}$，取得最小值.}
\step{故集合 $A\cap B$“长度”的取值范围是 $\left[\dfrac{1}{2},1\right]$.}
\solend

\fi
\item 已知集合 $M=\{x\mid 1\leqslant x\leqslant 10,\,x\in\mathbb{N}\}$，对它的非空子集 $A$，将 $A$ 中每个元素 $k$ 都乘以 $(-1)^k$ 再求和，如 $A=\{1,4,7\}$，可以求得和为 $(-1)^1\times 1+(-1)^4\times 4+(-1)^7\times 7=-4$，则对 $M$ 的所有非空子集，这些和的总和为 \blankline[2em].

\ans{$2560$}

\ifanswer
\solbegin
\step{$M=\{x\mid 1\leqslant x\leqslant 10,\,x\in\mathbb{N}\}=\{1,2,3,4,5,6,7,8,9,10\}$，$1$ 出现了 $2^9$ 次，其它元素一样，故集合 $M$ 的所有非空子集中，这些和的总和是 $2^9\times\left[(-1)^1\times 1+2\times(-1)^2+\dots+10\times(-1)^{10}\right]=2^9\times 5=2560$.}
\solend

\fi
\item 高一（13）班的同学发现对于一个有限集 $A$，一般都可以找到两个非空数集 $A_1$ 和 $A_2$，满足 $A_1\cap A_2=\varnothing$，且 $A_1\cup A_2=A$，记集合 $\{A_1,A_2\}$ 为 $A$ 的一个“划分集”. 若 $A$ 中有 $n\,(n\geqslant 2,\,n\in\mathbb{N}^*)$ 个元素，则不同的“划分集”共有 \blankline 个（用含 $n$ 的表达式填空）.

\ans{$2^{n-1}-1$}

\ifanswer
\solbegin
\step{集合中的每个元素都有属于 $A_1$ 或属于 $A_2$ 两种情况，故共有 $2^n$ 种情况，排除全都属于 $A_1$ 或 $A_2$ 的两种情况，考虑对称情况，故不同的“划分集”共有 $\dfrac{1}{2}(2^n-2)=2^{n-1}-1$ 个.}
\solend

\fi
\end{enumerate}

\bigsec{二、选择题}

\begin{enumerate}
\setcounter{enumi}{12}

\item 下面六个关系式：① $\varnothing\subsetneq\{a\}$；② $a\subsetneq\{a\}$；③ $\{a\}\subseteq\{a\}$；④ $\{a\}\in\{a,b\}$；⑤ $a\in\{a,b,c\}$；⑥ $\varnothing\in\{a,b\}$. 其中正确的是（\quad）.

A. ①②③④\qquad B. ③⑤⑥\qquad C. ①④⑤\qquad D. ①③⑤

\ans{D}

\item 下列关于集合运算的结论，错误的是（\quad）.

A. $\complement_U(A\cup B)=\complement_U A\cap\complement_U B$\qquad B. $A\cap(B\cap C)=(A\cap B)\cap C$

C. $A\cap(B\cup C)=(A\cap B)\cup(A\cap C)$\qquad D. $A\cup(B\cap C)=(A\cap B)\cup(A\cap C)$

\ans{D}

\item 已知集合 $A=\{1,2\}$，$B=\{a,a^2+1\}$，若 $A\cap B=\{1\}$，则 $a=$（\quad）.

A. $1$\qquad B. $-1$\qquad C. $0$\qquad D. $2$

\ans{C}

\ifanswer
\solbegin
\step{集合 $A=\{1,2\}$，$B=\{a,a^2+1\}$，$A\cap B=\{1\}$，所以 $\begin{cases}a=1 \\ a^2+1\neq 2\end{cases}$ 或 $\begin{cases}a\neq 2 \\ a^2+1=1\end{cases}$，解得 $a=0$.}
\step{故选 C.}
\solend

\fi

\item 如图，已知全集 $U=\R$，集合 $A=\{x\mid (2x-3)(x+1)\leqslant 0\}$，$B=\{x\mid x>0\}$，则图中阴影部分表示的集合为（\quad）.

A. $\{x\mid x\leqslant -1\}$\qquad B. $\{x\mid x<-1\}$\qquad C. $\left\{x\,\middle|\,x\leqslant 0\text{ 或 }x>\dfrac{3}{2}\right\}$\qquad D. $\left\{x\,\middle|\,x<0\text{ 或 }x>\dfrac{3}{2}\right\}$

\ans{B}

\ifanswer
\solbegin
\step{集合 $A=\left\{x\,\middle|\,-1\leqslant x\leqslant\dfrac{3}{2}\right\}$，而 $B=\{x\mid x>0\}$，则 $A\cup B=\{x\mid x\geqslant -1\}$，由韦恩图得，图中阴影部分表示的集合为 $\complement_U(A\cup B)=\{x\mid x<-1\}$.}
\step{故选 B.}
\solend

\fi
\end{enumerate}

\bigsec{三、解答题}

\begin{enumerate}
\setcounter{enumi}{16}

\item 设集合 $P=\{x\mid x^2-x-6<0\}$，$Q=\{x\mid 2a\leqslant x\leqslant a+3\}$.
\begin{enumerate}[label=(\arabic*)]
\item 若 $P\cup Q=P$，求实数 $a$ 的取值范围；
\item 若 $P\cap Q=\varnothing$，求实数 $a$ 的取值范围.
\end{enumerate}

\ans{（1）$(-1,0)\cup(3,+\infty)$；（2）$(-\infty,-5]\cup\left[\dfrac{3}{2},+\infty\right)$}

\ifanswer
\solbegin
\stepm{(1)}{$P=\{x\mid -2<x<3\}$，因为 $P\cup Q=P$，所以 $Q\subseteq P$.}
\step{当 $Q=\varnothing$ 时，$2a>a+3$，所以 $a>3$；}
\step{当 $Q\neq\varnothing$ 时，$a\leqslant 3$，因为 $Q\subseteq P$，所以 $\begin{cases}2a>-2 \\ a+3<3\end{cases}$，所以 $-1<a<0$.}
\step{所以实数 $a$ 的取值范围为 $(-1,0)\cup(3,+\infty)$.}
\stepm{(2)}{因为 $P\cap Q=\varnothing$，当 $Q=\varnothing$ 时，$2a>a+3$，所以 $a>3$；}
\step{当 $Q\neq\varnothing$ 时，$a\leqslant 3$，因为 $P\cap Q=\varnothing$，所以 $a+3\leqslant -2$ 或 $2a\geqslant 3$，所以 $a\leqslant -5$ 或 $\dfrac{3}{2}\leqslant a\leqslant 3$.}
\step{所以实数 $a$ 的取值范围为 $(-\infty,-5]\cup\left[\dfrac{3}{2},+\infty\right)$.}
\solend

\fi
\ansspace[6]
\item 实数构成的集合 $A$ 满足条件：若 $a\in A,\,a\neq 1$，则 $\dfrac{1}{1-a}\in A$.
\begin{enumerate}[label=(\arabic*)]
\item 求证：集合 $A$ 中不可能只含有一个元素；
\item 若 $2\in A$，则在 $A$ 中还有另外两个数，求出这两个数.
\end{enumerate}

\ans{（1）略；（2）$-1$ 和 $\dfrac{1}{2}$}

\ifanswer
\solbegin
\stepm{(1)}{假设集合 $A$ 中只含有一个元素，则 $a=\dfrac{1}{1-a}\Rightarrow a^2-a+1=0,\,\Delta=-3<0$，这与实数构成的集合 $A$ 矛盾，故集合 $A$ 中不可能只含有一个元素.}
\stepm{(2)}{$2\in A\Rightarrow\dfrac{1}{1-2}\in A\Rightarrow -1\in A\Rightarrow\dfrac{1}{1-(-1)}\in A\Rightarrow\dfrac{1}{2}\in A$，所以 $-1,\dfrac{1}{2}$ 为所求的两个数.}
\solend

\fi
\ansspace[5]
\item 设集合 $A=\{x\mid x^2+4x=0\}$，$B=\{x\mid x^2+2(a+1)x+a^2-1=0\}$，若 $B\subseteq A$，求实数 $a$ 的取值范围.

\ans{$\{a\mid a\leqslant -1\text{ 或 }a=1\}$}

\ifanswer
\solbegin
\step{$A=\{x\mid x^2+4x=0\}=\{-4,0\}$，$B=\{x\mid x^2+2(a+1)x+a^2-1=0\}$，$B\subseteq A$.}
\stepm{①}{若 $B=\varnothing$，则 $\Delta=4(a+1)^2-4(a^2-1)=8a+8<0$，则 $a<-1$；}
\stepm{②}{若 $B=\{0\}$ 或 $\{-4\}$，则 $\Delta=4(a+1)^2-4(a^2-1)=8a+8=0$，解得 $a=-1$，将 $a=-1$ 代入方程 $x^2+2(a+1)x+a^2-1=0$，得 $x^2=0$，得 $x=0$，即 $B=\{0\}$ 符合要求；}
\stepm{③}{若 $B=A=\{-4,0\}$，则 $\Delta=8a+8>0$，即 $a>-1$，即 $x^2+2(a+1)x+a^2-1=0$ 的两根分别为 $-4$、$0$，则有 $a^2-1=0$ 且 $-2(a+1)=-4$，则 $a=1$.}
\step{综上所述，实数 $a$ 的取值范围是 $\{a\mid a\leqslant -1\text{ 或 }a=1\}$.}
\solend

\fi
\ansspace[6]
\item 命题甲：关于 $x$ 的方程 $x^2+4mx+m=0$ 无实根；命题乙：关于 $x$ 的方程 $x^2-(m+1)x+m=0$ 有两个不相等的正根. 设命题甲、命题乙为真命题时实数 $m$ 的取值分别组成集合 $A$、$B$.
\begin{enumerate}[label=(\arabic*)]
\item 求集合 $A$、$B$；
\item 若命题甲、乙中有且仅有一个是真命题，求实数 $m$ 的取值范围.
\end{enumerate}

\ans{（1）$A=\left\{m\,\middle|\,0<m<\dfrac{1}{4}\right\}$，$B=\{m\mid m>0\text{ 且 }m\neq 1\}$；（2）$\left[\dfrac{1}{4},1\right)\cup(1,+\infty)$}

\ifanswer
\solbegin
\stepm{(1)}{命题甲：关于 $x$ 的方程 $x^2+4mx+m=0$ 无实根，所以 $16m^2-4m<0$，解得 $0<m<\dfrac{1}{4}$.}
\step{命题乙：关于 $x$ 的方程 $x^2-(m+1)x+m=0$ 有两个不相等的正根，所以 $\begin{cases}(m+1)^2-4m>0 \\ m+1>0 \\ m>0\end{cases}$，解得 $m>0$ 且 $m\neq 1$.}
\step{故 $A=\left\{m\,\middle|\,0<m<\dfrac{1}{4}\right\}$；}
\step{$B=\{m\mid m>0\text{ 且 }m\neq 1\}$.}
\stepm{(2)}{命题甲、乙中有且仅有一个是真命题，故 ① 甲真乙假，即 $\begin{cases}0<m<\dfrac{1}{4} \\ 0\geqslant m\text{ 或 }m=1\end{cases}$，故解集为 $\varnothing$；}
\stepm{②}{甲假乙真，即 $\begin{cases}0\geqslant m\text{ 或 }m\geqslant\dfrac{1}{4} \\ 0<m\text{ 且 }m\neq 1\end{cases}$，解得 $m\in\left[\dfrac{1}{4},1\right)\cup(1,+\infty)$.}
\step{由 ①② 得 $m\in\left[\dfrac{1}{4},1\right)\cup(1,+\infty)$.}
\solend

\fi
\ansspace[7]
\item 已知非空数集 $S$ 满足：对任意给定的 $x$、$y\in S$（$x$、$y$ 可以相同），有 $x+y\in S$ 且 $x-y\in S$.
\begin{enumerate}[label=(\arabic*)]
\item 哪个数一定是 $S$ 中的元素？说明理由；
\item 若 $S$ 是有限集，求 $S$；
\item 若 $S$ 中最小的正数为 5，求 $S$.
\end{enumerate}

\ans{（1）$0$；（2）$S=\{0\}$；（3）$S=\{n\mid n=5k,\,k\in\Z\}$}

\ifanswer
\solbegin
\stepm{(1)}{数字 $0$ 一定是 $S$ 中的元素.}
\step{若 $x=y=0$，$0\in S$，则 $x+y=0+0=0$，$x-y=0-0=0$，即 $x+y\in S$，$x-y\in S$.}
\stepm{(2)}{$S$ 为有限非空数集，假设 $S$ 中有非零元素 $m$，则有形如 $mk\,(k\in\mathbb{N},k>0)$ 的所有实数都是 $S$ 的元素，与 $S$ 是有限集矛盾，所以 $S=\{0\}$.}
\stepm{(3)}{$S$ 中最小的正数为 5，则 5 的正整数倍一定是 $S$ 中元素，又 $-x=0-x$，$0\in S$，故当 $x\in S$ 时，$-x\in S$，故所有形如 $5k$（$k$ 为整数）的数都是 $S$ 中的元素.}
\step{假设 $S$ 中有形如 $5k+r\,(r\in\{1,2,3,4\})$ 的元素，那么 $5k+r-5k=r\in S$，这与 $S$ 中最小正整数 5 矛盾.}
\step{综上，$S=\{n\mid n=5k,\,k\in\Z\}$.}
\solend

\fi
\ansspace[8]
\end{enumerate}

\bigsec{附加题}

\begin{enumerate}

\ansneed{7cm}
\item 对于正整数集合 $A=\{a_1,a_2,\dots,a_n\}\,(n\in\mathbb{N}^*,\,n\geqslant 3)$，如果去掉其中任意一个元素 $a_i\,(i=1,2,\dots,n)$ 之后，剩余的所有元素组成的集合都能分为两个交集为空集的集合，且这两个集合的所有元素之和相等，就称集合 $A$ 为“和谐集”.
\begin{enumerate}[label=(\arabic*)]
\item 判断集合 $\{1,2,3,4,5\}$ 是否为“和谐集”，并说明理由；
\item 求证：集合 $\{1,3,5,7,9,11,13\}$ 是“和谐集”；
\item 求证：若集合 $A$ 是“和谐集”，则集合 $A$ 中元素个数为奇数.
\end{enumerate}

\ans{（1）不是；（2）略；（3）略}

\ifanswer
\solbegin
\stepm{(1)}{对于集合 $\{1,2,3,4,5\}$，当去掉元素 2 时，剩余的所有元素之和为 13，不能分为两个交集为空集且这两个集合的所有元素之和相等的集合，所以集合 $\{1,2,3,4,5\}$ 不是“和谐集”.}
\stepm{(2)}{设 $A=\{1,3,5,7,9,11,13\}$，当去掉元素 1 时，有 $3+5+7+9=11+13$；}
\step{当去掉元素 3 时，有 $1+9+13=5+7+11$；}
\step{当去掉元素 5 时，有 $9+13=1+3+7+11$；}
\step{当去掉元素 7 时，有 $1+9+11=3+5+13$；}
\step{当去掉元素 9 时，有 $1+3+5+11=7+13$；}
\step{当去掉元素 11 时，有 $3+7+9=1+5+13$；}
\step{当去掉元素 13 时，有 $1+3+5+9=7+11$.}
\step{所以集合 $A=\{1,3,5,7,9,11,13\}$ 是“和谐集”.}
\stepm{(3)}{设“和谐集”$A=\{a_1,a_2,\dots,a_n\}$ 所有元素之和为 $M$.}
\step{由题意得 $M-a_i\,(i=1,2,\dots,n)$ 均为偶数，因此 $a_i\,(i=1,2,\dots,n)$ 的奇偶性相同.}
\step{（i）如果 $M$ 为奇数，则 $a_i\,(i=1,2,\dots,n)$ 也均为奇数，由于 $M=a_1+a_2+\dots+a_n$，所以 $n$ 为奇数.}
\step{（ii）如果 $M$ 为偶数，则 $a_i\,(i=1,2,\dots,n)$ 均为偶数，此时设 $a_i=2b_i$，则 $\{b_1,b_2,\dots,b_n\}$ 也是“和谐集”.}
\step{重复上述操作有限次，便可得各项均为奇数的“和谐集”.}
\step{此时各项之和也为奇数，集合 $A$ 中元素个数为奇数.}
\step{综上所述，集合 $A$ 中元素个数为奇数.}
\solend

\fi
\ansspace*[10]
\end{enumerate}

\end{document}